\begin{resumo}

    Ordens de produção de mostruários percorrem seis estações com roteiros distintos e hoje são sequenciadas de forma empírica. O problema é modelado como Job Shop Scheduling com roteiros variáveis e minimização do makespan. Comparam-se um modelo exato (CP-SAT), Busca Tabu, \textit{Simulated Annealing} e regras de despacho, em 50 instâncias de 10 a 30 ordens. O CP-SAT provou o ótimo em 37 das 50 instâncias; as metaheurísticas ficaram em média 2,6\% a 3,1\% acima do melhor valor conhecido, contra 12,2\% do sequenciamento atual simulado (FIFO). No caso-base, o ótimo reduz o makespan em 6,8\%.
    \vspace{\onelineskip}

    \noindent\textbf{Palavras-chave}: Job Shop. Makespan. Busca Tabu. CP-SAT. Mostruários.

\end{resumo}

% --- resumo em inglês ---
\begin{resumo}[Abstract]
\begin{otherlanguage*}{english}

    Production orders for sample displays go through six workstations with different routings and are currently sequenced empirically. The problem is modeled as a Job Shop Scheduling Problem with variable routings, minimizing the makespan. An exact CP-SAT model, Tabu Search, Simulated Annealing and dispatching rules are compared on 50 instances with 10 to 30 orders. CP-SAT proved optimality in 37 of the 50 instances; the metaheuristics were on average 2.6\% to 3.1\% above the best known value, against 12.2\% for the simulated current practice (FIFO). In the base case, the optimal schedule reduces the makespan by 6.8\%.

    \vspace{\onelineskip}

    \noindent\textbf{Keywords}: Job Shop. Makespan. Tabu Search. CP-SAT. Sample displays.
\end{otherlanguage*}
\end{resumo}
